\documentclass[12pt,a4paper]{article}

\usepackage[T1]{fontenc}
\usepackage[utf8]{inputenc}
\usepackage[margin=2.5cm, top=2cm]{geometry}
\usepackage{amsmath, amssymb, amsthm, mathtools}
\usepackage{fancyhdr}
\usepackage{enumitem}
\usepackage{hyperref}
\usepackage{microtype}
\usepackage{lmodern}
\usepackage{xcolor}

\definecolor{accent}{HTML}{2E4057}
\hypersetup{hidelinks}

% ── Header / footer ───────────────────────────────────────────────────────────
\pagestyle{fancy}
\fancyhf{}
\lhead{Your Name \\ Student ID: 12345678}
\chead{\textbf{Assignment N}}
\rhead{COSC/MTH XXXX \\ \today}
\cfoot{\thepage}
\renewcommand{\headrulewidth}{0.4pt}

% ── Problem / Solution environments ──────────────────────────────────────────
\newcounter{problem}
\setcounter{problem}{0}

\newenvironment{problem}[1][]{%
  \refstepcounter{problem}%
  \vspace{6pt}%
  \noindent\rule{\linewidth}{0.6pt}\\[-4pt]
  \noindent\textbf{\color{accent}Problem~\theproblem}%
  \ifx\relax#1\relax\else\ \textit{(#1)}\fi\par\smallskip
}{%
  \par\noindent\rule{\linewidth}{0.4pt}\vspace{4pt}%
}

\newenvironment{solution}{%
  \par\medskip\noindent\textit{Solution.}\enspace\ignorespaces
}{%
  \hfill$\blacksquare$\par\medskip
}

% ── Math shortcuts ────────────────────────────────────────────────────────────
\newcommand{\R}{\mathbb{R}}
\newcommand{\N}{\mathbb{N}}
\newcommand{\Z}{\mathbb{Z}}
\newcommand{\Q}{\mathbb{Q}}
\newcommand{\C}{\mathbb{C}}
\newcommand{\eps}{\varepsilon}
\DeclarePairedDelimiter{\abs}{\lvert}{\rvert}
\DeclarePairedDelimiter{\norm}{\lVert}{\rVert}

% ── Document ──────────────────────────────────────────────────────────────────
\title{\textbf{Assignment N}\\[4pt]\large COSC/MTH XXXX --- Course Name}
\author{Your Name \\ \small Student ID: 12345678}
\date{\today}

\begin{document}

\maketitle
\thispagestyle{fancy}

% ─────────────────────────────────────────────────────────────────────────────

\begin{problem}[5 marks]
  Prove that $\sqrt{2}$ is irrational.
\end{problem}

\begin{solution}
  Suppose for contradiction that $\sqrt{2} = p/q$ with $p, q \in \Z$, $q \neq 0$,
  and $\gcd(p, q) = 1$. Then $2q^2 = p^2$, so $p^2$ is even, hence $p$ is even.
  Write $p = 2k$; then $2q^2 = 4k^2$, giving $q^2 = 2k^2$, so $q$ is also even.
  This contradicts $\gcd(p,q) = 1$.
\end{solution}

% ─────────────────────────────────────────────────────────────────────────────

\begin{problem}[3 marks]
  Let $G = (V, E)$ be a finite simple graph with $n$ vertices.
  Show that $G$ has at least two vertices of the same degree.
\end{problem}

\begin{solution}
  The degree of any vertex lies in $\{0, 1, \ldots, n-1\}$, giving $n$ possible
  values for $n$ vertices. However, degrees $0$ and $n-1$ cannot both be achieved
  simultaneously (if one vertex is adjacent to all others, no vertex is isolated).
  Hence the degrees live in a set of at most $n-1$ distinct values, so by the
  Pigeonhole Principle at least two vertices share a degree.
\end{solution}

% ─────────────────────────────────────────────────────────────────────────────

\begin{problem}[4 marks]
  State and prove the AM-GM inequality for two non-negative reals.
\end{problem}

\begin{solution}
  \textbf{Claim:} For $a, b \geq 0$, $\dfrac{a+b}{2} \geq \sqrt{ab}$.

  \medskip
  Since $(\sqrt{a} - \sqrt{b})^2 \geq 0$, expanding gives
  $a - 2\sqrt{ab} + b \geq 0$, i.e.\ $a + b \geq 2\sqrt{ab}$.
  Dividing by $2$ yields the result. Equality holds iff $\sqrt{a} = \sqrt{b}$,
  i.e.\ $a = b$.
\end{solution}

\end{document}
